\documentclass[12pt,a4paper]{article}

\usepackage[utf8]{inputenc}
\usepackage[romanian]{babel}
\usepackage{hyperref}
\usepackage{graphicx}
\usepackage{listings}
\usepackage{xcolor}
\usepackage{geometry}
\usepackage{indentfirst}

\geometry{margin=2.5cm}

\definecolor{codegreen}{rgb}{0,0.6,0}
\definecolor{codegray}{rgb}{0.5,0.5,0.5}
\definecolor{codepurple}{rgb}{0.58,0,0.82}
\definecolor{backcolour}{rgb}{0.95,0.95,0.92}

\lstset{
    backgroundcolor=\color{backcolour},   
    commentstyle=\color{codegreen},
    keywordstyle=\color{magenta},
    numberstyle=\tiny\color{codegray},
    stringstyle=\color{codepurple},
    basicstyle=\ttfamily\small,
    breakatwhitespace=false,         
    breaklines=true,                 
    captionpos=b,                    
    keepspaces=true,                 
    numbers=left,                    
    numbersep=5pt,                  
    showspaces=false,                
    showstringspaces=false,
    showtabs=false,                  
    tabsize=2
}

\title{Documentație Proiect: Redy\\ \large Platformă Inteligentă de Agregare și Analiză a Știrilor}
\author{Andrei}
\date{Mai 2026}

\begin{document}

\maketitle
\newpage

\tableofcontents
\newpage

\section{Motivația alegerii temei și utilitatea aplicației}

În era digitală actuală, volumul de informații disponibil online este copleșitor. Utilizatorii se confruntă cu fenomenul de "infobezitate", unde filtrarea conținutului relevant devine o sarcină consumatoare de timp. Redy a fost conceput ca o soluție modernă pentru centralizarea, organizarea și analizarea automată a fluxurilor de știri (RSS/Atom).

\subsection{Utilitatea aplicației}
Aplicația Redy oferă următoarele avantaje:
\begin{itemize}
    \item \textbf{Centralizare:} Utilizatorul poate urmări zeci de surse de știri într-o singură interfață coerentă.
    \item \textbf{Analiză Automată:} Folosind Machine Learning, aplicația categorizează articolele și analizează sentimentul acestora (pozitiv, negativ, neutru).
    \item \textbf{Eficiență:} Extragerea automată a conținutului complet al articolelor elimină necesitatea de a naviga pe site-uri externe pline de reclame.
    \item \textbf{Personalizare:} Utilizatorii își pot organiza fluxurile în categorii și pot marca articolele favorite.
\end{itemize}

\section{Structura aplicației}

Proiectul Redy utilizează o arhitectură de tip microservicii/servicii distribuite, optimizată pentru performanță și scalabilitate.

\subsection{Organizarea conținutului informațional}
Sistemul este împărțit în patru componente principale:
\begin{enumerate}
    \item \textbf{Serverul Backend (Rust):} Nucleul aplicației, responsabil pentru API-ul REST, autentificare și coordonarea datelor.
    \item \textbf{Interfața Client (React):} O aplicație Single Page Application (SPA) modernă care oferă o experiență de utilizare fluidă.
    \item \textbf{ML Worker (Python):} Un serviciu specializat în procesarea limbajului natural (NLP).
    \item \textbf{Baza de Date și Mesageria:} PostgreSQL pentru stocare și NATS JetStream pentru comunicare asincronă.
\end{enumerate}

\subsection{Structuri de date utilizate}
Principalele entități stocate în baza de date includ:
\begin{itemize}
    \item \textbf{User:} Informații despre utilizatori și setări.
    \item \textbf{Feed:} Sursa de știri (URL, nume).
    \item \textbf{Article:} Datele brute ale articolelor (titlu, link, conținut HTML).
    \item \textbf{ArticleData:} Metadate generate de procesele de ML (categoria, scorul de sentiment).
    \item \textbf{Category:} Grupări de fluxuri definite de utilizator.
\end{itemize}

\newpage

\section{Tehnologii de bază utilizate}

Această secțiune prezintă limbajele și instrumentele utilizate, explicând de ce acestea reprezintă standardul actual în industrie și cum se raportează ele la conceptele studiate în curricula școlară.

\subsection{Limbajul de programare Rust: Siguranță și Performanță}

Deși nu este studiat în liceu, Rust este considerat "succesorul spiritual" al limbajelor C și C++. În timp ce C++ oferă programatorului control total asupra memoriei, acesta vine cu riscul unor erori grave (precum accesarea unei zone de memorie inexistente), care pot duce la prăbușirea programului sau la vulnerabilități de securitate.

\textbf{De ce Rust?}
\begin{itemize}
    \item \textbf{Compilatorul ca "Profesor":} În Rust, compilatorul este mult mai strict decât în C++. Acesta verifică nu doar sintaxa, ci și \textit{logica gestionării memoriei}. Dacă un programator încearcă să facă o operațiune nesigură, programul pur și simplu nu se va compila, oferind mesaje de eroare extrem de detaliate care explică "de ce" nu este corect.
    \item \textbf{Viteză fără compromisuri:} Rust este un limbaj \textbf{compilat} (ca și C++), ceea ce înseamnă că este transformat direct în cod mașină pe care procesorul îl înțelege nativ. Acest lucru îl face ideal pentru servere care trebuie să proceseze mii de știri simultan cu un consum minim de curent și resurse.
    \item \textbf{Fără "Garbage Collector":} Spre deosebire de limbaje precum Java sau C#, Rust nu are nevoie de un proces secundar care să "curețe" memoria în timpul rulării, ceea ce elimină pauzele neașteptate în execuție.
\end{itemize}

\begin{lstlisting}[language=Rust, caption=Exemplu de rigoare in Rust]
fn main() {
    let mesaj = String::from("Salut");
    let copie = mesaj; // Proprietatea asupra textului se muta la 'copie'
    
    // Urmatoarea linie ar opri compilarea in Rust:
    // println!("{}", mesaj); 
    // Motivul: 'mesaj' nu mai detine datele, prevenind utilizarea dubla a memoriei.
}
\end{lstlisting}

\newpage

\subsection{Limbajul de programare Python: Limbajul Inteligenței Artificiale}

Dacă Rust este despre viteză și precizie, Python este despre \textbf{simplitate și flexibilitate}. Este adesea primul limbaj învățat de studenți datorită sintaxei sale care seamănă foarte mult cu limba engleză.

\textbf{Rolul Python în proiect:}
\begin{itemize}
    \item \textbf{Limbaj Interpretat:} Spre deosebire de Rust, Python este un limbaj \textbf{interpretat}, ceea ce înseamnă că instrucțiunile sunt citite și executate rând cu rând de către un program special (interpretat). Deși este mai lent la calcule matematice brute, simplitatea sa îl face ideal pentru a coordona biblioteci complexe de calcul.
    \item \textbf{Standardul în Inteligență Artificială:} Aproape toate descoperirile recente în domeniul AI (precum ChatGPT sau recunoașterea facială) sunt dezvoltate în Python. Acest lucru se datorează numărului imens de "pachete" (biblioteci de cod gata scrise) disponibile gratuit.
    \item \textbf{Sintaxă "Curată":} Python elimină nevoia de acolade și punct și virgulă, folosind \textbf{indentarea} (spațiile de la începutul rândului) pentru a structura codul, ceea ce forțează programatorul să scrie cod ordonat și ușor de citit.
\end{itemize}

\begin{lstlisting}[language=Python, caption=Simplitatea Python in procesarea textului]
# In Python, operatiile complexe devin usor de citit
text = "Redy este o platforma inteligenta"
cuvinte = text.split() # Imparte textul in cuvinte

for cuvânt in cuvinte:
    print(f"Am gasit cuvantul: {cuvânt}")
\end{lstlisting}

\newpage

\subsection{Biblioteca React: Construirea Interfețelor Moderne}

React nu este un limbaj de programare propriu-zis, ci o bibliotecă construită peste \textbf{JavaScript} (limbajul care rulează în browser). React a revoluționat modul în care construim site-uri web, trecând de la pagini statice la aplicații web dinamice.

\textbf{Concepte explicate:}
\begin{itemize}
    \item \textbf{Gândirea în Componente:} În loc să scriem o pagină web ca pe un document lung, o vedem ca pe un set de piese de "LEGO". Avem o componentă pentru butoane, una pentru cardul de știri și una pentru meniul de sus. Acestea pot fi refolosite oriunde în aplicație.
    \item \textbf{Actualizare Automată:} În dezvoltarea web clasică, dacă voiai să schimbi un preț pe ecran, trebuia să "cauți" manual elementul și să-l modifici. În React, programatorul doar descrie cum arată interfața în funcție de date, iar React se ocupă automat de actualizarea ecranului atunci când datele se schimbă (de exemplu, când apare o știre nouă).
    \item \textbf{Interactivitate Fluidă:} Utilizarea React permite paginii să se actualizeze instant, fără a fi nevoie de o reîncărcare completă (refresh), oferind o experiență similară cu o aplicație de telefon mobil.
\end{itemize}

\begin{lstlisting}[language=JavaScript, caption=Exemplu de componenta React]
function ButonApreciere() {
  const [aprecieri, setAprecieri] = React.useState(0);

  return (
    <button onClick={() => setAprecieri(aprecieri + 1)}>
      Voturi: {aprecieri}
    </button>
  );
}
\end{lstlisting}

\newpage

\subsection{Comunicarea între servicii via NATS}
NATS JetStream permite componentelor să comunice fără a fi strâns cuplate, crescând reziliența sistemului.

\textbf{Mecanismul Pub/Sub:}
\begin{itemize}
    \item \textbf{Publish:} Backend-ul emite un eveniment (ex: "Articol Nou").
    \item \textbf{Subscribe:} Worker-ul ML recepționează evenimentul și îl procesează.
\end{itemize}

\begin{lstlisting}[language=Rust, caption=Emiterea unui eveniment (Publish)]
// Backend-ul "arunca" sarcina in coada si trece mai departe
js.publish("tasks.ml", payload.into()).await?;
\end{lstlisting}

\newpage

\section{Detalii tehnice de implementare}

Această secțiune descrie în detaliu arhitectura sistemului și soluțiile tehnice adoptate pentru a asigura performanța, securitatea și scalabilitatea platformei Redy.

\subsection{Backend-ul în Rust: Performanță și Securitate}

Serverul central este pilonul de rezistență al aplicației, fiind construit pe o stivă tehnologică modernă care prioritizează securitatea datelor și execuția asincronă.

\textbf{1. Arhitectura API-ului (Axum):}
Utilizăm framework-ul \texttt{axum}, care se bazează pe biblioteca \texttt{tokio} pentru runtime-ul asincron. Un aspect tehnic avansat este utilizarea \textbf{Extractorilor}, care permit validarea și preluarea automată a datelor din cererile HTTP (starea bazei de date, autentificarea utilizatorului etc.).

\textbf{2. Securitatea și Autentificarea (OIDC):}
Pentru a asigura un nivel înalt de securitate, platforma utilizează protocolul \textbf{OpenID Connect (OIDC)}, permițând utilizatorilor să se conecteze în siguranță folosind furnizori externi (ex: Google). Implementarea include măsuri avansate precum codurile \textbf{PKCE (Proof Key for Code Exchange)} pentru a preveni interceptarea codurilor de autorizare.

\begin{lstlisting}[language=Rust, caption=Generarea URL-ului de autentificare securizat]
pub fn get_auth_url(&self) -> (String, CsrfToken, Nonce, PkceCodeVerifier) {
    let client = CoreClient::from_provider_metadata(/*...*/)
        .set_redirect_uri(self.redirect_url.clone());

    // Generam elemente de securitate unice pentru fiecare sesiune
    let (pkce_challenge, pkce_verifier) = PkceCodeChallenge::new_random_sha256();
    let (auth_url, csrf_token, nonce) = client
        .authorize_url(
            CoreAuthenticationFlow::AuthorizationCode,
            CsrfToken::new_random,
            Nonce::new_random,
        )
        .add_scope(Scope::new("email".to_string()))
        .set_pkce_challenge(pkce_challenge)
        .url();

    (auth_url.to_string(), csrf_token, nonce, pkce_verifier)
}
\end{lstlisting}

\textbf{3. Gestiunea Datelor (SeaORM):}
Interacțiunea cu PostgreSQL este realizată printr-un ORM asincron (\texttt{SeaORM}). Acesta ne permite să utilizăm \textbf{ActiveModels} pentru operații de tip CRUD sigure și migrații de bază de date versionate. Un exemplu de logică implementată este procesul de "Find or Create", unde utilizatorul este identificat după email sau creat automat la prima logare.

\newpage

\subsection{Procesarea ML în Python: Inteligență Asincronă}

Componenta de Machine Learning este proiectată să funcționeze ca un set de workeri independenți, conectați la backend prin NATS JetStream.

\textbf{1. Managementul Modelelor:}
Pentru a optimiza consumul de memorie (modelele pot ocupa peste 2GB RAM), am implementat un sistem de încărcare de tip \textbf{Lazy Loading/Singleton}. Modelele sunt încărcate o singură dată la pornirea worker-ului și sunt partajate între toate sarcinile de procesare primite.

\textbf{2. Procesul de Clasificare (Zero-Shot):}
Spre deosebire de clasificarea tradițională care necesită re-antrenare pentru categorii noi, Redy utilizează tehnica \textit{Zero-Shot Classification}. Acest lucru permite sistemului să categorizeze știri în orice domeniu nou definit de utilizator, doar prin furnizarea numelui categoriei ca etichetă (label).

\begin{lstlisting}[language=Python, caption=Logica de clasificare dinamica]
def handle_categorize(article_uuid: uuid.UUID):
    # Preluam toate categoriile disponibile din sistem
    categories = repository.get_all_human_category_names()
    text = utils.extract_text_from_html(article.html_content)
    
    # Modelul calculeaza probabilitatea pentru fiecare eticheta
    result = model.classify(text, categories)
    top_category = result['labels'][0] # Cea mai probabila categorie
    
    repository.update_article_category(article_uuid, top_category)
\end{lstlisting}

\textbf{3. Reziliența prin Mesagerie:}
Utilizarea \texttt{nats-py} cu confirmări manuale (\texttt{manual\_ack=True}) asigură că nicio știre nu rămâne neprocesată. Dacă un model ML eșuează din cauza lipsei de resurse, mesajul este retrimis automat către un alt worker disponibil.

\newpage

\subsection{Frontend-ul în React: Interfață Reactivă și Tipizată}

Clientul web este o aplicație modernă construită cu \texttt{TypeScript}, oferind o experiență similară unei aplicații native.

\textbf{1. Managementul Stării (TanStack Query):}
În loc să gestionăm manual \texttt{useEffect} și \texttt{loading states}, utilizăm \textbf{TanStack Query}. Această bibliotecă oferă un mecanism robust de \textit{caching}, \textit{background refetching} și \textit{synchronization} între client și server. 

\begin{lstlisting}[language=JavaScript, caption=Sincronizarea automata a datelor cu TanStack Query]
export function FeedArticleList({ feedUuid }) {
    // Cererea este cache-uita si re-executata doar daca datele devin "stale"
    const { data, isLoading } = $api.useQuery("get", "/articles", {
        params: { query: { feed_uuid: feedUuid } }
    });

    if (isLoading) return <LoadingSkeleton />;
    
    return (
        <div className="grid gap-4">
            {data.map(item => <ArticleCard key={item.id} {...item} />)}
        </div>
    );
}
\end{lstlisting}

\textbf{2. Integrarea cu API-ul și TypeScript:}
Utilizăm un generator de tipuri care transformă definiția OpenAPI a backend-ului în interfețe TypeScript. Acest lucru garantează că orice schimbare în structura datelor de pe server va fi detectată la compilarea frontend-ului, eliminând erorile de tip "undefined field".

\textbf{3. Componente UI (Shadcn/UI):}
Interfața este construită folosind un sistem de componente bazat pe \textbf{Radix UI} și \textbf{Tailwind CSS}, asigurând accesibilitate maximă și o estetică modernă (Dark/Light mode) fără a compromite performanța de randare.

\newpage

\section{Resurse hard și soft necesare}

\subsection{Resurse Soft}
\begin{itemize}
    \item \textbf{Limbaje:} Rust 1.75+, Python 3.10+, Node.js 18+.
    \item \textbf{Baze de date:} PostgreSQL 14+.
    \item \textbf{Infrastructură:} NATS Server, Docker.
\end{itemize}

\subsection{Resurse Hard}
\begin{itemize}
    \item \textbf{Procesor:} Minim 2 nuclee (recomandat 4 pentru ML).
    \item \textbf{RAM:} Minim 4GB (modelele ML pot fi intensive).
    \item \textbf{Stocare:} 10GB pentru baza de date și containere.
\end{itemize}

\section{Modalități de utilizare}

\subsection{Instalare și Pornire}
Aplicația folosește Docker Compose pentru o pornire rapidă:
\begin{enumerate}
    \item Clonarea proiectului.
    \item Comanda: \texttt{docker-compose up -d}.
    \item Accesare: \texttt{http://localhost:3000}.
\end{enumerate}

\subsection{Fluxul Utilizatorului}
Utilizatorul se autentifică, adaugă un flux RSS de interes (ex: știri tehnice), iar sistemul începe automat fetch-ul și analiza. Articolele apar în dashboard cu etichete de sentiment și categorie.

\section{Posibilități de dezvoltare}

\begin{itemize}
    \item \textbf{Sistem de Recomandări:} Bazat pe istoricul de citire.
    \item \textbf{Aplicație Mobilă:} Dezvoltată în React Native.
    \item \textbf{Suport Multilingv:} Extinderea modelelor ML pentru limba română.
\end{itemize}

\end{document}
